\documentclass[11pt]{article}
\usepackage[a4paper,margin=22mm]{geometry}
\usepackage{fontspec}
\setmainfont{DejaVu Serif}
\setsansfont{DejaVu Sans}
\usepackage{microtype}
\usepackage{xcolor}
\usepackage{hyperref}
\usepackage{enumitem}
\usepackage{parskip}
\definecolor{Accent}{HTML}{971D32}
\definecolor{Muted}{HTML}{666666}
\hypersetup{colorlinks=true,urlcolor=Accent,linkcolor=Accent}
\setlist{leftmargin=1.6em,itemsep=0.3em,topsep=0.35em}
\setcounter{tocdepth}{2}
\renewcommand{\contentsname}{Contents}
\newcommand{\sectionrule}{\par\vspace{-0.5em}\noindent\textcolor{Accent}{\rule{\linewidth}{0.6pt}}\par}
\begin{document}

\begin{center}
{\sffamily\bfseries\color{Accent} TOEFL WORD OF THE DAY\par}
\vspace{8mm}
{\fontsize{42}{48}\selectfont\bfseries pioneer\par}
\vspace{2mm}
{\Large\sffamily\color{Accent} /ˌpaɪəˈnɪr/\par}
\vspace{2mm}
{\small\sffamily Local audio phrase: pioneer\par}
{\small\sffamily Audio file: audios/pioneer.mp3\par}
\vspace{2mm}
{\large\sffamily\color{Muted} countable noun and transitive verb\par}
\vspace{5mm}
{\sffamily an early innovator \textbullet\ an early explorer or settler \textbullet\ to develop something first\par}
\vspace{6mm}
{\large A pioneer is among the first to enter a place or develop an important idea, method, or field.\par}
\vspace{6mm}
\begin{minipage}{0.84\textwidth}
\itshape “The scientist was a pioneer in the study of human memory.”
\end{minipage}\par
\vspace{10mm}
\textbf{Date:} July 26th 2026\quad\textbullet\quad \textbf{Level:} B1--mid-B2 TOEFL
\vfill
Created and compiled by \textbf{Amir Kooshky}\par
\vspace{2mm}
\href{https://t.me/KooshkyTOEFL}{Telegram Channel}\quad\textbullet\quad
\href{https://www.instagram.com/kooshkytoefl}{Instagram}\quad\textbullet\quad
\href{https://t.me/KooshkyTOEFL\_pv}{Personal Telegram}
\end{center}
\newpage
\tableofcontents
\newpage

\section{Meanings and usage}
\sectionrule
\subsection{Meaning 1: Early innovator}
\textbf{Part of speech:} countable noun

\textbf{IPA:} /ˌpaɪəˈnɪr/

\textbf{Pronunciation phrase:} pioneer

\textbf{Local audio file:} audios/pioneer.mp3

\textbf{Register:} neutral to formal; common in academic, historical, and professional contexts

\textbf{Definition:} a person who is among the first to develop a new idea, method, technology, or area of study

\textbf{In simpler words:} one of the first people to develop something important

\textbf{Typical pattern:} a pioneer in + field or a pioneer of + development

\subsubsection{Natural examples}
\begin{enumerate}
  \item The scientist was a pioneer in the study of human memory.
  \item She is widely regarded as a pioneer of modern environmental economics.
  \item Early pioneers in computer science created many of the basic ideas used today.
  \item The medical pioneer introduced a safer way to perform the operation.
  \item As a pioneer in online education, the university began offering digital courses decades ago.
  \item The exhibition honors pioneers whose inventions transformed communication.
\end{enumerate}

\subsubsection{Nuance and implication}
\textbf{Central nuance:} The word emphasizes being early and helping open a new path for later people.

\textbf{Usual implication:} It often suggests originality, influence, courage, or major contribution.

\textbf{Compared with founder:} A founder creates an organization or institution, while a pioneer is among the first to develop or advance an idea, method, or field. One person can be both.

\subsubsection{Grammar notes}
\textbf{How to use it:} Use a pioneer for one person and pioneers for more than one.

\textbf{What can follow it:} It is often followed by in and the name of a field, or by of and the development that the person helped begin.

\textbf{Common preposition:} Use pioneer in for an area of activity, as in a pioneer in genetics. Use pioneer of for a movement, method, or development, as in a pioneer of modern surgery.

\textbf{Noun use:} This is a countable noun, so use a, an, or the with the singular form.

\subsubsection{Useful patterns}
\textbf{Pattern:} a pioneer in + field

\textbf{Use:} Use this to identify the area in which someone was among the first.

\textbf{Example:} She became a pioneer in renewable-energy research.

\textbf{Pattern:} a pioneer of + method, movement, or development

\textbf{Use:} Use this to name the specific thing that someone helped introduce.

\textbf{Example:} He was a pioneer of evidence-based teaching.

\textbf{Pattern:} one of the pioneers of + development

\textbf{Use:} Use this when several people helped begin or advance something.

\textbf{Example:} The engineer was one of the pioneers of satellite communication.

\subsubsection{Common wrong usage}
\paragraph{Correction 1}
\textbf{Wrong:} She is pioneer in robotics.

\textbf{Correct:} She is a pioneer in robotics.

\textbf{Why:} The singular countable noun normally needs a, an, or the.

\paragraph{Correction 2}
\textbf{Wrong:} He was a pioneer on artificial intelligence.

\textbf{Correct:} He was a pioneer in artificial intelligence.

\textbf{Why:} Use in before a field or area of activity.

\subsection{Meaning 2: Early explorer or settler}
\textbf{Part of speech:} countable noun

\textbf{IPA:} /ˌpaɪəˈnɪr/

\textbf{Pronunciation phrase:} pioneer

\textbf{Local audio file:} audios/pioneer.mp3

\textbf{Register:} neutral; especially common in historical writing

\textbf{Definition:} one of the first people to explore, travel to, or settle in a particular area

\textbf{In simpler words:} an early explorer or settler

\textbf{Typical pattern:} early pioneers or pioneers who settled + place

\subsubsection{Natural examples}
\begin{enumerate}
  \item The pioneers crossed the mountains in search of new farmland.
  \item Life was difficult for the early pioneers who settled the region.
  \item The museum displays tools used by pioneers during the westward expansion.
  \item Many pioneers traveled long distances without reliable maps.
  \item The valley's first pioneers built homes near the river.
  \item Historians study how pioneers interacted with the Indigenous communities already living there.
\end{enumerate}

\subsubsection{Nuance and implication}
\textbf{Central nuance:} This sense refers to people who entered or settled an area early from the viewpoint of the group being discussed.

\textbf{Usual implication:} Historical use may describe hardship and exploration, but it should not suggest that the land was empty or uninhabited.

\subsubsection{Grammar notes}
\textbf{How to use it:} Use a pioneer for one person and pioneers for a group.

\textbf{What can follow it:} A place, route, or historical period is often added later in the sentence to show where or when the pioneers traveled or settled.

\textbf{Noun use:} This is a countable noun. The plural pioneers is especially common when discussing a historical group.

\subsubsection{Useful patterns}
\textbf{Pattern:} early pioneers

\textbf{Use:} Use this for people who were among the first members of a group to enter or settle an area.

\textbf{Example:} Early pioneers established farms along the coast.

\textbf{Pattern:} pioneers who settled + place

\textbf{Use:} Use this to identify a group by the place where they settled.

\textbf{Example:} The book follows pioneers who settled the northern plains.

\subsubsection{Common wrong usage}
\paragraph{Correction 1}
\textbf{Wrong:} The pioneer people crossed the desert.

\textbf{Correct:} The pioneers crossed the desert.

\textbf{Why:} Use pioneers as the noun. Pioneer people is not the usual expression.

\subsection{Meaning 3: Develop or introduce first}
\textbf{Part of speech:} transitive verb

\textbf{IPA:} /ˌpaɪəˈnɪr/

\textbf{Pronunciation phrase:} pioneer

\textbf{Local audio file:} audios/pioneer.mp3

\textbf{Register:} neutral to formal; frequent in academic, scientific, historical, and business writing

\textbf{Definition:} to be among the first to develop, introduce, or use a new idea, method, technology, or activity

\textbf{In simpler words:} to develop or use something before most others

\textbf{Typical pattern:} pioneer + method, technology, approach, or the use of something

\subsubsection{Natural examples}
\begin{enumerate}
  \item The research team pioneered a new method for detecting the disease.
  \item The company pioneered the use of recycled materials in its products.
  \item These educators pioneered an approach that later spread worldwide.
  \item The laboratory pioneered several techniques for studying ancient DNA.
  \item Local farmers pioneered water-saving methods during a severe drought.
  \item The organization helped pioneer community-based health care in rural areas.
\end{enumerate}

\subsubsection{Nuance and implication}
\textbf{Central nuance:} The verb stresses that a person or group opened a new path by creating or using something early.

\textbf{Usual implication:} It usually suggests that the innovation influenced later work.

\subsubsection{Grammar notes}
\textbf{How to use it:} This verb normally needs the new thing after it.

\textbf{What can follow it:} It is commonly followed by words such as method, technique, approach, technology, process, or the use of something.

\textbf{Position in a sentence:} The person, team, or organization that did the early work normally comes before the verb.

\subsubsection{Useful patterns}
\textbf{Pattern:} pioneer + a method, technique, or approach

\textbf{Use:} Use this when someone develops a new way of doing something.

\textbf{Example:} The clinic pioneered a technique for repairing damaged tissue.

\textbf{Pattern:} pioneer the use of + noun

\textbf{Use:} Use this when someone is among the first to use a particular material, tool, or practice.

\textbf{Example:} The architects pioneered the use of natural ventilation in large buildings.

\textbf{Pattern:} help pioneer + noun

\textbf{Use:} Use this when someone contributed to early development but did not work alone.

\textbf{Example:} The professor helped pioneer a new model of language instruction.

\subsubsection{Common wrong usage}
\paragraph{Correction 1}
\textbf{Wrong:} The scientists pioneered in a new treatment.

\textbf{Correct:} The scientists pioneered a new treatment.

\textbf{Why:} In this meaning, pioneer normally takes the new idea, method, or technology directly after it.

\paragraph{Correction 2}
\textbf{Wrong:} The company pioneered to use electric vehicles.

\textbf{Correct:} The company pioneered the use of electric vehicles.

\textbf{Why:} Use pioneer the use of something, not pioneer to use something.

\section{Word family}
\sectionrule
\subsection{pioneering}
\textbf{Part of speech:} adjective

\textbf{IPA:} /ˌpaɪəˈnɪrɪŋ/

\textbf{Definition:} introducing new ideas, methods, or technologies before most others

\textbf{Relationship to the headword:} This adjective describes work or people that pioneer something.

\textbf{Grammar pattern:} pioneering + work, research, study, method, or technology

\textbf{Usage note:} It is usually positive and emphasizes originality and early influence.

\subsubsection{Examples}
\begin{enumerate}
  \item Her pioneering research changed the treatment of infectious diseases.
  \item The award recognized the team's pioneering work in solar engineering.
\end{enumerate}

\section{Common collocations}
\sectionrule
\subsection{a pioneer in the field of}
\textbf{Applies to:} Early innovator

\textbf{Meaning:} one of the first important people working in a particular field

\textbf{Pattern:} a pioneer in the field of + subject

\textbf{Example:} She was a pioneer in the field of developmental psychology.

\subsection{a pioneer of modern}
\textbf{Applies to:} Early innovator

\textbf{Meaning:} a person who helped begin or shape a modern development

\textbf{Pattern:} a pioneer of modern + field or practice

\textbf{Example:} The architect is considered a pioneer of modern urban planning.

\subsection{early pioneer}
\textbf{Applies to:} Early innovator

\textbf{Meaning:} a person involved at a very early stage

\textbf{Example:} She was an early pioneer of educational television.

\subsection{pioneer a new method}
\textbf{Applies to:} Develop or introduce first

\textbf{Meaning:} develop a method before most others

\textbf{Pattern:} pioneer + a new method

\textbf{Example:} The researchers pioneered a new method of analyzing climate data.

\subsection{pioneer the use of}
\textbf{Applies to:} Develop or introduce first

\textbf{Meaning:} be among the first to use something

\textbf{Pattern:} pioneer the use of + noun

\textbf{Example:} The hospital pioneered the use of robots in delicate operations.

\subsection{pioneering research}
\textbf{Applies to:} Develop or introduce first

\textbf{Meaning:} highly original research that opens a new area

\textbf{Example:} Her pioneering research provided the foundation for later discoveries.

\textbf{Usage note:} Pioneering is an adjective in this collocation.

\subsection{pioneering work}
\textbf{Applies to:} Develop or introduce first

\textbf{Meaning:} original early work that strongly influences later developments

\textbf{Example:} Their pioneering work made low-cost water purification possible.

\textbf{Usage note:} Pioneering is an adjective in this collocation.

\section{Synonyms and antonyms}
\sectionrule
\subsection{Close synonyms}
\subsubsection{innovator}
\textbf{Part of speech:} countable noun

\textbf{Applies to:} Early innovator

\textbf{Shared meaning:} Both words can describe a person who develops something new.

\textbf{Difference:} An innovator introduces new ideas or methods. A pioneer also emphasizes being among the earliest people to open a new area or path.

\textbf{Example:} The award honors innovators who create practical solutions to environmental problems.

\subsubsection{trailblazer}
\textbf{Part of speech:} countable noun

\textbf{Applies to:} Early innovator

\textbf{Shared meaning:} Both words describe an influential person who enters or develops an area early.

\textbf{Difference:} Trailblazer is more vivid and often more informal. It strongly suggests creating a path that others later follow.

\textbf{Example:} She was a trailblazer for women in aerospace engineering.

\subsubsection{settler}
\textbf{Part of speech:} countable noun

\textbf{Applies to:} Early explorer or settler

\textbf{Shared meaning:} Both can refer to people who establish homes in a new place.

\textbf{Difference:} A settler is anyone who moves to live permanently in a place. A pioneer is specifically among the earliest members of the group to do so.

\textbf{Example:} The settlers built permanent homes near the harbor.

\subsubsection{develop}
\textbf{Part of speech:} transitive verb

\textbf{Applies to:} Develop or introduce first

\textbf{Shared meaning:} Both can describe creating a new method, technology, or approach.

\textbf{Difference:} Develop means to create or improve something. Pioneer adds the idea of doing so before most other people.

\textbf{Example:} The team developed a more accurate forecasting model.

\section{Practice exercises}
\sectionrule
\subsection{Word family choice}
Choose the best member of the word family for each blank.

\begin{enumerate}
  \item The scientist received an international award for her \_\_\_\_\_ research on gene therapy.
  \item The engineer was a \_\_\_\_\_ in the development of low-cost solar cells.
\end{enumerate}

\subsection{Correct or incorrect?}
Decide whether each sentence is correct. Correct the inaccurate ones.

\begin{enumerate}
  \item The laboratory pioneered a rapid test for the infection.
  \item She was pioneer of modern nursing.
  \item The company pioneered to use biodegradable packaging.
\end{enumerate}

\subsection{Collocation practice}
Complete each sentence with a natural collocation.

\begin{enumerate}
  \item The institute helped \_\_\_\_\_ the use of satellite images in agricultural research.
  \item Her \_\_\_\_\_ work in public health influenced policies around the world.
\end{enumerate}

\subsection{Complete answer key}
\subsubsection{Word family choice}
\begin{enumerate}
  \item \textbf{pioneering.} The scientist received an international award for her pioneering research on gene therapy. \emph{Use the adjective pioneering before the noun research.}
  \item \textbf{pioneer.} The engineer was a pioneer in the development of low-cost solar cells. \emph{Use the countable noun pioneer after a.}
\end{enumerate}

\subsubsection{Correct or incorrect?}
\begin{enumerate}
  \item \textbf{Correct} \emph{The verb is correctly followed directly by the new thing that the laboratory developed.}
  \item \textbf{Incorrect}. She was a pioneer of modern nursing. \emph{The singular countable noun needs a.}
  \item \textbf{Incorrect}. The company pioneered the use of biodegradable packaging. \emph{Use pioneer the use of something.}
\end{enumerate}

\subsubsection{Collocation practice}
\begin{enumerate}
  \item \textbf{pioneer.} The institute helped pioneer the use of satellite images in agricultural research. \emph{Help pioneer the use of is a natural pattern for contributing to an early innovation.}
  \item \textbf{pioneering.} Her pioneering work in public health influenced policies around the world. \emph{Pioneering work is a common collocation for original, influential early work.}
\end{enumerate}

\vfill
\begin{center}
\small\color{Muted} Created and compiled by \textbf{Amir Kooshky}\\
\href{https://t.me/KooshkyTOEFL}{Telegram Channel}\quad\textbullet\quad
\href{https://www.instagram.com/kooshkytoefl}{Instagram}\quad\textbullet\quad
\href{https://t.me/KooshkyTOEFL\_pv}{Personal Telegram}
\end{center}
\end{document}
